\section{The Exact Period-Eight Spectral Edge}

The rational number in Section 4 is convenient for proving strict finite
counterexamples, but it is not the true band edge. We now prove Theorem C.

\subsection{The endpoint roots}

At \(c=2\), equation (4.5) becomes

\begin{lstlisting}[language={}]
P(y,2)=y^4-16y^3+76y^2-96y+16.                                 (5.1)
\end{lstlisting}

Translate \(y=x+4\). Then

\begin{lstlisting}[language={}]
P(x+4,2)=x^4-20x^2+80.                                          (5.2)
\end{lstlisting}

Putting \(w=x^2\) gives \(w^2-20w+80=0\), whose roots are
\(10\pm 2\sqrt{5}\). The four real roots of (5.1), in increasing order, are

\begin{lstlisting}[language={}]
4-sqrt(10+2sqrt(5)),
4-sqrt(10-2sqrt(5)),
4+sqrt(10-2sqrt(5)),
4+sqrt(10+2sqrt(5)).                                             (5.3)
\end{lstlisting}

Thus the largest endpoint root is

\begin{lstlisting}[language={}]
eta=4+sqrt(10+2sqrt(5)).                                        (5.4)
\end{lstlisting}

\subsection{Sharp positivity and uniqueness}

Set

\begin{lstlisting}[language={}]
s=sqrt(10+2sqrt(5)),       eta=4+s,
u=y-eta,                   t=2-c.
\end{lstlisting}

Substitution in (4.5) yields the identity

\begin{lstlisting}[language={}]
P(eta+u,2-t)
 =u^4+4s u^3+2u^2t+(40+12sqrt(5))u^2
  +4s ut+8sqrt(5)s u+t^2+(4sqrt(5)-3)t.                         (5.5)
\end{lstlisting}

All coefficients are positive. Hence for \(u,t\ge 0\), the right side is
nonnegative, and it vanishes only at \(u=t=0\). Since every unit-circle fiber
has \(c\le 2\), equation (5.5) excludes every squared eigenvalue above \(\eta\) and
permits equality only when \(c=2\).

At \(c=2\), equation (5.1) has \(\eta\) as a root, so \(\pm \sqrt{\eta}\) are eigenvalues
of \(H(1)\). Therefore

\begin{lstlisting}[language={}]
sup_(|z|=1) rho(H(z))^2=eta.                                    (5.6)
\end{lstlisting}

On the unit circle, \(z+z^{-1}=2\) if and only if \(z=1\). This proves both the
value and the uniqueness assertion in Theorem C. \(\square\)

For later use we also record that \(\eta\) has minimal polynomial

\begin{lstlisting}[language={}]
Y^4-16Y^3+76Y^2-96Y+16                                         (5.7)
\end{lstlisting}

Indeed, translation \(Y=X+4\) reduces (5.7) to
\(X^4-20X^2+80\). By Gauss's lemma, a rational factorization may be taken into
monic integral quadratics. Vanishing of the cubic and linear coefficients
reduces it either to \((X^2+a)(X^2+b)\), where
\(a+b=-20\) and \(ab=80\), or to
\((X^2+rX+s)(X^2-rX+s)\), where \(s^2=80\). The first case would make \(a,b\)
roots of \(T^2+20T+80\), whose discriminant 80 is not a square; the second has
no rational \(s\). Thus (5.7) is irreducible over \(Q\). The number \(\eta\) lies in
\((1951/250,1561/200)\); in particular \(\eta<8\).

\subsection{Monotonicity of the top band}

Let \(r(c)\) be the largest real root of \(P(y,c)\). At \(c=-2\),

\begin{lstlisting}[language={}]
P(y,-2)=(y^2-12y+34)(y^2-4y+2),
r(-2)=6+sqrt(2)=:y_0.                                           (5.8)
\end{lstlisting}

For \(-2<c\le 2\), substitution gives

\begin{lstlisting}[language={}]
P(y_0,c)=(c+2)(c+5-8sqrt(2))<0.                                (5.9)
\end{lstlisting}

Since \(P(y,c)\) tends to positive infinity with \(y\), (5.9) implies
\(r(c)>y_0\) for \(c>-2\).

Moreover,

\begin{lstlisting}[language={}]
P_c(y,c)=2(c-c_0(y)),       c_0(y)=y^2-8y+13/2.                 (5.10)
\end{lstlisting}

The function \(c_0\) is increasing for \(y\ge y_0>4\), and

\begin{lstlisting}[language={}]
c_0(y_0)=-7/2+4sqrt(2)>2.
\end{lstlisting}

Thus \(P_c(y,c)<0\) throughout \(y\ge y_0\) and \(c\le 2\). If
\(-2\le c_1<c_2\le 2\) and \(y_1=r(c_1)\), then
\(P(y_1,c_2)<P(y_1,c_1)=0\); hence \(P(.,c_2)\) has a root above \(y_1\).
Therefore \(r(c)\) is strictly increasing on \([-2,2]\).

\subsection{Finite holonomies}

For \(\alpha=+1\), the finite set \(z^L=1\) contains \(z=1\). By (5.6),

\begin{lstlisting}[language={}]
rho(A_(8L,+1))^2=eta.                                           (5.11)
\end{lstlisting}

For \(\alpha=-1\), the admissible parameters are
\(z_k=exp((2k+1)\pi i/L)\). Their largest \(c\) value is \(2\operatorname{cos}(\pi/L)\), which is
strictly less than two. The monotonicity just proved gives

\begin{lstlisting}[language={}]
rho(A_(8L,-1))^2=r(2cos(pi/L))<eta.                             (5.12)
\end{lstlisting}

As \(L\) tends to infinity, \(2\operatorname{cos}(\pi/L)\) tends to two, and continuity of the
eigenvalues of the Hermitian fiber gives convergence of (5.12) to \(\eta\).
Thus both holonomy sectors have the same infinite-volume edge, but only the
positive sector attains it at each finite size.
